% Source: https://artofproblemsolving.com/community/c3h3590075_varied_number_theory
% Stable identity: aops-topic-3590075-q02
Write $\dfrac{n}{40-n}=k^2$, $k\ge0$. Then $n=40-\dfrac{40}{k^2+1}$. Thus $k^2+1$ must divide $40$. Since $k^2+1\le40$, check $k=0,1,2,3,4,5,6$; the divisors occur at $k=0,1,2,3$, yielding $n=0,20,32,36$. The answer is $\boxed4$.
